\chapter{The last-level cache datapath}\label{ch:llc-datapath}

\whatyoufind{The shipped LLC drawn as a datapath, and the same drawing with one request's path
highlighted four times: a read hit, a write hit, a read miss into a free way, and a miss that
evicts. These five figures are the concrete instance of everything \cref{ch:cache-controller}
describes, wired up the way the four-core system actually configures it, and they are where
the latencies the simulator reports come from.}

\section{The datapath}

\begin{figure}[htbp]\centering\includegraphics[width=\textwidth]{llc_architecture.pdf}
\caption[The shipped LLC as a datapath]{The shipped LLC as a datapath: \code{CacheController\_End2End}
with \code{CacheDataHandler\_COTS}, 32\,KB, 2-way, 64\,B lines, the \code{MESI\_LLC} FSM and
one timed data port. Bus interfaces on top, control on the left, the single array in the
middle read by an untimed bits path and a timed data port on the right, MSHR and write-back
buffer beside it, the memory bus below.}\label{fig:llc-architecture}\end{figure}

\Cref{fig:llc-architecture} follows the model rather than a textbook cache. The address splits
into tag, set and offset, and the set indexes \emph{one} array: in Octopus a line carries its
coherence bits and its data together, so there is no separate tag array. That array is read two
ways, and the difference matters for every latency the simulator reports. Reading a line's bits
--- the tag compare, the state the FSM needs --- is \textbf{untimed} and claims nothing. Reading
or writing its \textbf{data} goes through a single port, which demand reads, eviction reads and
fill or write-back writes contend for through the round-robin arbiter, and which is then busy
for \code{A\_LLC} cycles. A line that currently sits in the MSHR or the write-back buffer is read
without waiting for that port, because it is not in the array yet --- the dotted bypass --- but
it still restarts the busy timer, which is why such accesses cut in ahead of one already in
progress.

Control is the dashed amber layer: the line state and the decoded event enter the coherence
FSM, whose next state is written back into the line and whose actions become port claims,
buffer allocations and messages. A miss allocates in the MSHR and leaves on the memory bus; a
victim moves to the write-back buffer and frees its way at once; and because the cache is
inclusive, the FSM's \act{IssueInv} on an eviction invalidates the copy in whichever L1 holds the
line --- the red path, and the reason a task whose working set fits in its own private cache can
still be disturbed by other cores (\cref{ch:tasks-and-demo}).

\section{A read hit}

\figwide{llc_hit_path.pdf}{A read or write hit --- a \code{GetS} or \code{GetM} that finds the
line in a stable state --- from the request's arrival on the bus to its data leaving on the
response lane, in ten numbered steps.}{fig:llc-hit}

Steps 4--6 --- tag compare, FSM, sequencer --- take no simulated time. What a hit pays for is the
wait in the processing queue (1--3), the wait for the port (7), and then the \code{A\_LLC}-cycle
access itself (8--9). The figure draws the case where the LLC serves the data: the line is valid
in the array and no L1 owns it (\state{I} or \state{S} in \code{MESI\_LLC.csv}). The third case,
a hit on a line an L1 holds in \state{EorM}, does not use the port at all: the owner answers the
\code{GetS} on the bus with \act{Data2Both}, and the LLC only saves that copy
(\state{S\_d} $\rightarrow$ \act{SaveData} $\rightarrow$ \state{S}); on a \code{GetM} the owner
hands the line straight to the requester and the LLC merely rewrites the owner bits.

\section{A write hit}

\figwide{llc_write_hit.pdf}{The write hit on its own. The path is the read hit's; what is new
is the branch marked 5b --- the FSM's \code{SetOwner} written back into the line's bits over the
dashed control wire, untimed --- and the plain-data response.}{fig:llc-write-hit}

A write hit differs from a read hit in what the FSM does, not in the path: the \code{GetM} gets
plain data where a \code{GetS} gets exclusive data, and the FSM also rewrites the owner bits.
Nothing else changes --- in particular the LLC sends \textbf{no} invalidation: on the snooping
bus every sharer sees the \code{GetM} itself and drops its copy
(\state{S} + \code{Other\_GetM} $\rightarrow$ \state{I}), so \act{IssueInv} is reserved for
evictions.

\section{A read miss into a free way}

\figwide{llc_read_miss.pdf}{A read miss in the simplest case: the set still has a free way, so
nothing is evicted. The request's own path in teal; the array write of the fill, off the
critical path, dashed.}{fig:llc-read-miss}

The request takes the same first six steps and then diverges at the tag compare: the FSM's
\act{GetData} allocates an MSHR entry, the read leaves on \code{bus[1]}, and the line waits in the
MSHR --- not the array --- until DRAM answers. When the fill is processed the response goes out
on TX response at once, carrying the data from the fill message; the write into the free way
--- which contends for the port like any other array access and frees the MSHR --- runs in
parallel, off the requester's critical path. Because no victim is chosen, the write-back buffer
and the red inclusion path are never touched. A write miss takes the identical route with a
\code{GetM} in place of the \code{GetS}.

\section{A miss that evicts}

\figwide{llc_miss_evict.pdf}{A miss into a \emph{full} set: the requested line in teal, the
evicted victim in amber. The victim is chosen when the fill is written, and its exit is a
coherence transaction of its own.}{fig:llc-miss-evict}

Two things in \cref{fig:llc-miss-evict} are easy to get wrong from a textbook. First,
\textbf{the victim is chosen late}: the miss allocates an MSHR and sends the read while the full
set is left alone, and only when the fill's array write runs does the data handler pick the LRU
line among the requester's allowed ways, move it --- bits and data --- into the write-back
buffer, and put the fill in its way (\code{moveLine2WB} inside \code{updateLineData}). The
requester's data leaves on TX response right then; everything the victim costs comes afterwards
and lands on \emph{other} traffic. Second, \textbf{the victim's exit is a coherence transaction},
not a buffer drain: a \code{Replacement} is raised for it and goes through the queue and the FSM
like any request, \act{IssueInv} puts an \code{INV} on the service lane, every L1 that shares the
line drops it, the LLC receives its own \code{INV} back, and only then does \act{WriteBack} read
the line --- from the buffer, over the dotted bypass --- and send it to memory. There is no dirty
bit in the model, so every victim is written back; the code marks the spot with a \code{ToDo}.
A \code{GetM} miss evicts identically. A victim an L1 \emph{owns} adds one leg: the owner answers
the \code{INV} with its data, and that copy is what reaches memory.

Seen from the cores instead of from inside the LLC, that eviction is the whole story of
inter-core interference through an inclusive LLC; \cref{ch:tasks-and-demo} draws it as a
sequence and \cref{ch:request-end-to-end} shows what it costs the victim's next request.

\begin{wherebox}
\begin{itemize}[nosep,leftmargin=1.2em]
\item \file{src/CacheControllers/CacheController_End2End.cpp} --- the LLC controller
\item \file{src/CacheDataHandler_COTS.cpp} --- \code{updateLineData}, \code{moveLine2WB}, the port arbiter and the bypass
\item \file{Protocols_FSM/MESI_LLC.csv} --- the FSM (\cref{app:fsm-tables})
\item \file{src/Protocols/LLCMESIProtocol.cpp} --- its handler
\item \file{src/CacheControllers/CacheController.cpp} --- \code{sendInvalidationMessage}, \code{checkReplacements}
\item \file{configuration/SystemConfigurations/MultiCoreSystem.csv} --- the \code{llc\_controller.*} keys that size it
\end{itemize}
\end{wherebox}
